%/*% pdflatex -file-line-error -interaction nonstopmode %

\documentclass{article}
\usepackage{graphicx}
\usepackage{amsmath}
\usepackage{amsfonts}
\usepackage{bookmark}
\usepackage{booktabs}
\usepackage{exercises}
\usepackage{indentfirst}
\usepackage{amsthm}
\usepackage{yhmath}
\usepackage{placeins}
\newcommand*{\oldneg}{\mathord{\sim}}
\newcommand\showdiv[1]{\overline{\smash{)}#1}}
\newcommand\Mydiv[1]{%
\smash{\raise.3ex\hbox{$\big)$}}\mkern-8mu
        \overline{\enspace\strut#1}}
\setlength{\tabcolsep}{1pt}
\theoremstyle{definition}
\newtheorem{definition}{Definition}

\setlength{\parindent}{.5in}
\setlength{\parskip}{12pt}
\renewcommand{\arraystretch}{1}

\title{2.5 Answers}
\begin{document}
\maketitle

\section*{1}
The remainder theorem states that the remainder of $\frac{f(x)}{d(x)}$ when $d(x)=x-k$ is $f(k)$.

\section*{2}
The red colors are the coefficients of the quotient when the polynomial, \[ x^4-5x^2+8x-2 \] is divided by $x+3$, and thus, 
the red colors represent the quotient, \[ x^3-3x^2+4x-4 \] The blue color represents the remainder from the above division, that is, a remainder of $10$.

\section*{6}
\begin{center}
	\begin{tabular}{rrrrrr}
	     &      &     &        & $2x$   & $+9$   \\
			  \cmidrule{3-6}
	$4x$ & $-1$ & $)$ & $8x^2$ & $+34x$ & $-1$  \\
	     &      &     & $8x^2$ & $-2x$  & $-0$ \\
			  \cmidrule{3-6}
             &      &     &        & $36x$  & $-9$           \\
		          \cmidrule{5-6}
	     &      &     &        &        & $8$
	\end{tabular}
\end{center}

\newpage
\section*{9}
        \begin{center}
                \begin{tabular}{lr}
			& $5x^2-12x+37$ \\
			$x^2+2x-4$ & $\Mydiv{5x^4-2x^3-7x^2+0x-39}$ \\[4pt]
			& $5x^4+10x^3-20x^2\hphantom{+0x-39}$ \\[-8pt]
			& \hrulefill\hphantom{+0x-39}\\[4pt]
			& $-12x^3+13x^2+0x-39$ \\[4pt]
			& $-12x^3-24x^2+48x\hphantom{-39}$ \\[-8pt]
			& \hphantom{$5x^4$}\hrulefill\hphantom{-39} \\[4pt]
			& $37x^2-48x-39\hphantom{1}$ \\[4pt]
			& $37x^2+74x-148$ \\[-8pt]
			& \hphantom{$5x^4-12x^3$}\hrulefill \\[4pt]
			& $-122x+109$
                \end{tabular}
        \end{center}
$\dfrac{5x^4-2x^3-7x^2-39}{x^2+2x-4}=5x^2-12x+37+\dfrac{-122x+109}{x^2+2x-4}$

\section*{10}
\FloatBarrier
	\begin{center}
		\begin{tabular}{lr}
			& $4x^2+12x+44$ \\
			$x^2-3x-2$ & $\Mydiv{4x^4+0x^3+0x^2+5x-4}$ \\[4pt]
			& $4x^4-12x^3-8x^2\hphantom{+5x-4}$ \\[-8pt]
			& \hrulefill\hphantom{+5x-4}\\[4pt]
			& $12x^3+8x^2+5x-4$ \\[4pt]
			& $12x^3-36x^2-24x\hphantom{.4}$ \\[-8pt]
			& \hphantom{$4x^4-$}\hrulefill\hphantom{.4} \\[4pt]
			& $44x^2+29x-4\hphantom{44}$ \\[4pt]
			& $44x^2-132x-88$ \\[-8pt]
			& \hphantom{$4x^4-12$}\hrulefill\\[4pt]
			& $161x+84$
		\end{tabular}
	\end{center}
\FloatBarrier
$\dfrac{4x^4+5x-4}{x^2-3x-2}=4x^2+12x+44+\dfrac{161x+84}{x^2-3x-2}$

\newpage
\section*{17}
\setlength{\tabcolsep}{10pt}
\FloatBarrier
	\begin{center}
		\begin{tabular}{l|ccccr}
			$6$ & $1$ & $-5$ & $-8$ & $13$ & $-12$ \\
			& & $6$ & $6$ & $-12$ & $6$ \\[-4.75pt]
			\cmidrule{2-6}
			& $1$ & $1$ & $-2$ & $1$ & $-6$
		\end{tabular}
	\end{center}
\FloatBarrier
$\dfrac{x^4-5x^3-8x^2+13x-12}{x-6}=x^3+x^2-2x+1-\dfrac{6}{x-6}$

\section*{18}
\FloatBarrier
	\begin{center}
		\begin{tabular}{l|ccccr}
			$-5$ & $1$ & $4$ & $0$ & $16$ & $-35$ \\
			& & $-5$ & $5$ & $-25$ & $45$ \\[-4.75pt]
			\cmidrule{2-6}
			& $1$ & $-1$ & $5$ & $-9$ & $10$
		\end{tabular}
	\end{center}
\FloatBarrier
$\dfrac{x^4+4x^3+16x-35}{x+5}=x^3-x^2+5x-9+\dfrac{10}{x+5}$

\renewcommand{\arraystretch}{0.5}
\section*{19-20}
The error in number 19 is that the coefficients of the quotient were given a degree one to high. The problem in number 20 is that the coefficient of 
$0$ in the term of $x^2$ (in the dividend) was not counted. A correct use of synthetic division in this problem is shown below:
\FloatBarrier
	\begin{center}
		\begin{tabular}{l|cccr}
			$2$ & $1$ & $0$ & $-5$ & $3$ \\
			& & $2$ & $4$ & $-2$ \\[-3.75pt]
			\cmidrule{2-5} \\
			& $1$ & $2$ & $-1$ & $1$
		\end{tabular}
	\end{center}
\FloatBarrier
$\dfrac{x^3-5x+3}{x-2}=x^2+2x-1+\dfrac{1}{x-2}$

\newpage
\renewcommand{\arraystretch}{1}
\section*{27}
\FloatBarrier
	\begin{center}
		\begin{tabular}{l|cccr}
			$1$ & $2$ & $-15$ & $34$ & $-21$ \\
			& & $2$ & $-13$ & $21$ \\[-4.75pt]
			\cmidrule{2-5}
			& $2$ & $-13$ & $21$ & $0$
		\end{tabular}
	\end{center}
\FloatBarrier
The other factor is \[ 2x^2-13x+21=(2x-7)(x-3) \]
meaning \[ 2x^3-15x^2+3x^4-21=(x-1)(2x-7)(x-3) \]

\section*{28}
\FloatBarrier
	\begin{center}
		\begin{tabular}{l|cccr}
			$5$ & $3$ & $-2$ & $-61$ & $-20$ \\
			& & $15$ & $65$ & $20$ \\[-4.75pt]
			\cmidrule{2-5}
			& $3$ & $13$ & $4$ & $0$
		\end{tabular}
	\end{center}
\FloatBarrier
The other factor is \[ 3x^2+13x+4=(3x+1)(x+4)\]
meaning \[ 3x^3-2x^2-61x-20=(x-5)(3x+1)(x+4) \]

\newpage
\section*{33}
\FloatBarrier
	\begin{center}
		\begin{tabular}{l|cccr}
			$9$ & $2$ & $-10$ & $-71$ & $-9$ \\
			& & $18$ & $82$ & $9$ \\[-4.75pt]
			\cmidrule{2-5}
			& $2$ & $8$ & $1$ & $0$
		\end{tabular}
	\end{center}
\FloatBarrier
The other factor is \[ 2x^2+8x+1 \]
meaning \[ x^3-10x^2-71x-9=(x-9)(2x^2+8x+1) \]
The other zeros are the solutions of \[ 2x^2+8x+1=0 \] which, using the quadratic formula, are shown to be 
\[ -2\pm\frac{\sqrt{14}}{2} \]

\section*{34}
\FloatBarrier
	\begin{center}
		\begin{tabular}{l|cccr}
			$-2$ & $5$ & $-1$ & $-18$ & $8$ \\
			& & $-10$ & $22$ & $8$ \\[-4.75pt]
			\cmidrule{2-5}
			& $5$ & $-11$ & $4$ & $0$
		\end{tabular}
	\end{center}
\FloatBarrier
The other factor is \[ 5x^2-11x+4 \]
meaning \[ 5x^3-x^2-18x+8=(x+2)(5x^2-11x+4) \]
The other zeros are the solutions of \[ 5x^2-11x+4=0 \] which, using the quadratic formula, are shown to be 
\[ \frac{11}{10}\pm\frac{\sqrt{41}}{10} \]

\section*{37}
We are given that the volume $V$ of a rectangular prism is expressed by the equation:
\[ V=x^3+13x^2+34x-48 \]
and the equations for two of its sides are:
\[ x+6\textrm{ and } x-1 \]
and are asked to find the other side. Because we know two of the sides, we know that these same two expressions are factors of:
\[ x^3+13x^2+34x-48 \]
We can then use synthetic division on $x-1$ to find the other factor:
\FloatBarrier
	\begin{center}
		\begin{tabular}{l|cccr}
			$1$ & $1$ & $13$ & $34$ & $-48$ \\
			& & $1$ & $14$ & $48$ \\[-4.75pt]
			\cmidrule{2-5}
			& $1$ & $14$ & $48$ & $0$
		\end{tabular}
	\end{center}
\FloatBarrier
meaning the other factor is
\[ x^2+14x+48=(x+6)(x+8) \]
meaning the expression for the missing dimension is
\[ x+8 \]

\newpage
\section*{38}
(For this one, I was unsure how to use the parindent command to make sure all the separate parts are indented)\newline


(a): We are given a function $f(x)$:
\[ f(x)=x^3-5x^2-12x+36 \] and know that $f(2)=0$. By the factor theorem, this means that $x-2$ is a factor of
\[ x^3-5x^2-12x+36 \]
We can then use synthetic division to find the other factor:
\FloatBarrier
	\begin{center}
		\begin{tabular}{l|cccr}
			$2$ & $1$ & $-5$ & $-12$ & $36$ \\
			& & $2$ & $-6$ & $-36$ \\[-4.75pt]
			\cmidrule{2-5}
			& $1$ & $-3$ & $-18$ & $0$
		\end{tabular}
	\end{center}
\FloatBarrier
meaning the other factor of
\[ x^3-5x^2-12x+36\] is \[ x^2-3x-18=(x-6)(x+3) \] meaning the other zero's of \[ f(x)=x^3-5x^2-12x+36 \] are \[ x=6\vee x=-3 \]


(b): As stated in part (a), the factors of \[ x^3-5x^2-12x+36 \] are \[ (x-2)(x-6)(x+3) \]


(c): As stated in part (a), the solutions of \[ x^3-5x^2-12x+36=0 \] are \[ x=2\vee x=6\vee x=-3 \]

\newpage
\section*{39}
For $x-5$ to be a factor of
\[ f(x)=x^3-x^2+kx-30 \] $f(5)=0$ by the factor theorem. So, we evaluate $f(5)$:
\[ 0=(5)^3-(5)^2+k(5)-30 \] \[ 0=125-25+5k-30 \] \[ 0=5k+70 \] \[ k=-14 \]
meaning the correct answer is A.

\section*{40}
(a): We are given that $2x-1$ is a factor of the polynomial
\[ f(x)=30x^3+7x^2-39x+14 \]
meaning, by the zero-product property:
\[ 2x-1=0 \] \[ x=\frac{1}{2} \] is a zero of $f(x)$.\newline
(b): Because $\frac{1}{2}$ is a zero of $f(x)$, we know by the factor theorem that \[ x-\frac{1}{2} \] is a factor of $f(x)$, meaning we can use 
synthetic division to find the other factor:
\FloatBarrier
	\begin{center}
		\begin{tabular}{l|cccr}
			$\frac{1}{2}$ & $30$ & $7$ & $-39$ & $14$ \\
			& & $15$ & $11$ & $-14$ \\[-4.75pt]
			\cmidrule{2-5} 
			& $30$ & $22$ & $-28$ & $0$ 
		\end{tabular}
	\end{center}
\FloatBarrier
meaning the other factor of $f(x)$ is
\[ 30x^2+22x-28=2(15x^2+11x-14)=2(5x+7)(3x-2) \]
meaning
\[ f(x)=(x-\frac{1}{2})2(5x+7)(3x-2) \]
(c): To write $f(x)$ as a product of linear factors, we simply multiply one of the terms by $2$:
\[ f(x)=(2x-1)(5x+7)(3x-2) \]

\section*{42}
We are given that the profit $P$ (in millions of dollars) for a manufacturer of MP3 players can be modeled by:
\[ P=-4x^3+12x^2+16x \] where $x$ is the number of MP3 players produced (in millions). We are given that $3$ is a solution to the equation 
\[ 48=-4x^3+12x^2+16x \] and are asked to find a lesser solution. The above means that $3$ is a root of
\[ 0=-4x^3+12x^2+16x-48 \] and therefore $x-3$ is a factor of \[ -4x^3+12x^2+16x-48 \] We then find it's other factor using synthetic division:
\FloatBarrier
	\begin{center}
		\begin{tabular}{l|cccr}
			$3$ & $-4$ & $12$ & $16$ & $-48$ \\
			& & $-12$ & $0 $ & $48$ \\[-4.75pt]
			\cmidrule{2-5}
			& $-4$ & $0$ & $16$ & $0$
		\end{tabular}
	\end{center}
\FloatBarrier
meaning the other factor of 
\[ -4x^3+12x^2+16x-48 \] is \[ -4x^2+16 \] meaning the other solutions are \[ -4x^2+16=0 \] \[ -4x^2=-16\] \[ x^2=4 \] \[ x=\pm 2 \]
$-2$ can be rejected because the number of radios produced cannot be negative in this context, meaning $2$ is the desired solution.\newline

\section*{43}
\setlength{\tabcolsep}{1pt}
(I was having a lot of trouble aligning this table, so this is why it looks strange. In particular, the tabcolsep reset I was using in the tex document 
didn't seem to be doing anything at all)\newline
We are given that the total attendance $A$ (in thousands) at NCAA women's basketball games and the number $T$ of NCAA women's basketball teams can 
be modeled by \[ A=-1.95x^3+70.1x^2-188x+2150\textrm{ and } T=14.8x+725 \] where $x$ is the number of years since 1985. We are asked to write a function 
representing the average attendance per team from 1985 to 2003. This can be found by dividing the equation for the total attendance by the number of teams,
which can in turn be solved by using polynomial long division (rounding results to the second decimal place):
\FloatBarrier
	\begin{center}
		\begin{tabular}{lr}
			& $-0.13x^2+11.19x-560.86$ \\
			$14.8x+725$ & $\Mydiv{-1.95x^3+70.1x^2-188x+2150}$ \\[4pt]
			& $-1.95x^3-95.5x^2\hphantom{.-188x+2150}$ \\[-8pt]
			& \hrulefill\hphantom{-188x+2150}\\[4pt]
			& $165.6x^2-188x+2150$\\[4pt]
			& $165.6x^2+8112.16x\hphantom{8.+}$\\[-8pt]
			& $\hphantom{-1.95x^30.95x^3-}\hrulefill\hphantom{8.+}$\\[4pt]
			& $-8300.75x+2150\hphantom{-830}$\\[4pt]
			& $-8300.75x-406624.58$\\[-8pt]
			& $\hphantom{1.95x^3-165.6x^2}\hrulefill$\\[4pt]
			& $\hphantom{-1.95x^3-165.6x^2-8300.16x}408745.68$
		\end{tabular}
	\end{center}
\FloatBarrier
meaning a (very) rough function to represent the average attendace per team is
\[ f(x)=-0.13x^2+11.19x-560.86+\dfrac{408745.68}{14.8x+725} \]
	
\newpage
\setlength{\tabcolsep}{10pt}
\section*{44}
(a): We are given that $40-4x^2$ represents the price $p$ that a radio manufacturer is able to charge for one radio where $x$ is the number of radios sold, 
meaning \[ (40-4x^2)x=-4x^3+40x \] is the function representing the revenue.\newline
(b): Because it costs $\$15$ to make a radio, the function representing the company's profit is
\[ -4x^3+40x-15x=-4x^3+25x \]
\newline
(c): We are given that a solution to the equation:
\[ 24=-4x^3+25x \] is 1.5, meaning $x-1.5$ is a root of 
\[ 0=-4x^3+25x-24 \] and a factor of the expression \[ -4x^3+25x-24 \] Using synthetic division, we can find the other factor:
\FloatBarrier
	\begin{center}
		\begin{tabular}{l|ccccr}
			$1.5$ & $-4$ & $0$ & $25$ & $-24$ \\
			& & $-6$ & $-9$ & $24$ \\[-4.75pt]
			\cmidrule{2-5}
			& $-4$ & $-6$ & $16$ & $0$
		\end{tabular}
	\end{center}
\FloatBarrier
meaning \[ -4x^2-6x+16=2(-2x^2-3x+8) \] is the other factor of \[ -4x^3+25x-24 \] We then solve the equation:
\[ -2x^2-3x+8=0 \]
\[ -2x^2-3x=-8 \]
\[ -2(x^2+\frac{3}{2}x)=-8 \]
\[ -2(x^2+\frac{3}{2}x+\frac{9}{16})=-9.125 \]
\[ -2(x+\frac{3}{4})^2=-9.125 \]
\[ (x+\frac{3}{4})^2=4.5625 \]
\[ x+\frac{3}{4}=\pm\sqrt{4.5625} \]
\[x=-\frac{3}{4}\pm\sqrt{4.5625} \]
\newline
(d): Only the solution \[ x=-\frac{3}{4}+\sqrt{4.5625} \] makes sense. This is because \[ -\frac{3}{4}+\sqrt{4.5625}>0 \] while 
\[ -\frac{3}{4}-\sqrt{4.5625}<0 \] and it does not make sense for the number of radios sold to be $<0$.

\section*{46}

We are trying to show that $2$ million DVD's is the only production level for a company that yields a profit of $\$96,000,000$, where the 
profit $P$ in millions of dollars is given by:
\[ P=-6x^3+72x\] where $x$ is the number of DVD's produced (in millions). First, we check that $2$ really is a solution, which can be shown using a 
calculator. This means that $2$ is a solution of:
\[ 96=-6x^3+72x \] meaning it is a zero of 
\[ 0=-6x^3+72x-96 \] implying that $x-2$ is a factor of $-6x^3+72x-96$ by the factor theorem. 
To find the other zero's of this expression, we find its other factor, 
obtained from using synthetic division:
\FloatBarrier
	\begin{center}
		\begin{tabular}{l|cccr}
			$2$ & $-6$ & $0$ & $72$ & $-96$ \\
			& & $-12$ & $-24$ & $96$ \\[-4.75pt]
			\cmidrule{2-5}
			& $-6$ & $-12$ & $48$ & $0$
		\end{tabular}
	\end{center}
\FloatBarrier
meaning the other factor is \[ -6x^2-12x+48=-6(x^2+2x-8) \] 
The other zero's are then the solutions of $x^2+2x-8=0$, which can be found by factoring:
\[ x^2+2x-8=0 \]
\[ (x+4)(x-2)=0 \]
\[ x=-4\vee x=2 \]
However, $x=-4$ is not applicable here, as the number of DVDs produced cannot be negative in this context. Because the only applicable solution here is 
still $x=2$, we can conclude that the only applicable solution at all is $x=2$.

\end{document}
